\section{绪论}
\label{sec:introduction}

\subsection{研究背景与意义}

货币政策通过政策工具的实际操作改变金融条件，也通过公开沟通改变市场对未来政策路径和经济前景的判断。中国人民银行按季发布的《中国货币政策执行报告》（下称 MPR）同时记录经济运行、政策执行和下一阶段政策思路，为在统一发布节奏下研究书面沟通提供了文本样本。报告中较积极的措辞可能伴随较好的增长预期。市场也可能据此修正对进一步宽松的预期。前者主要涉及预期现金流，后者可能通过贴现率发挥作用，两种机制对价格的作用方向未必一致。

行业层面的问题是，市场收到同一份报告后，各行业的相对收益能否呈现稳定的排序。即使股票市场整体在发布后上涨，各行业若以近似幅度反应，市场平均收益的变化也难以转化为行业配置价值。相反，行业的融资结构、现金流久期和政策暴露不同，可能使相同信息在行业间产生差异。因而，检验行业收益的增量可预测性，需要显式区分发布附近的异常收益、跨行业的相对收益，以及投资者在信号可得之后能够实现的收益。

\subsection{文献综述}

\subsubsection{MPR 文本与股票市场}
姜富伟等从报告构造文本情绪、相似度和可读性指标，发现情绪改善与发布后市场正向反应相关，相关证据集中在货币政策指引部分\cite{jiangHuHuang2021}。许魁结合央行报告与分析师研报考察信息传递，讨论了分析师解读在央行文本与股票市场之间的作用\cite{xuKui2023}。林雨欣分别构造央行沟通指数与实际行动指数，在非金融上市公司样本中比较行业分组的收益与波动反应\cite{linYuxin2025}。这些研究为区分政策立场和市场反应提供了基础，本文另行检验报告全文情绪在申万一级行业指数上的样本外预测表现。

\subsubsection{央行沟通的宏观与风险传导}
王辉等将报告情绪与市场解读并列，发现市场接收程度和前瞻性措辞与银行系统性风险有关\cite{wangZhuHu2023}。蒋海等将前瞻性文本情绪与银行流动性囤积联系起来\cite{jiangSunLi2025}。王熙、郑梦圆将文本语调纳入结构向量自回归，以区分货币政策冲击及其伴随的信息效应\cite{wangZheng2021}。纪钧尧等使用 FinBERT 比较多经济体央行文本，发现文本情绪与汇率波动的联系比与汇率水平的联系更为明显\cite{jiYouChen2026}。上述结果也提示，沟通的内容、受众和结果变量会改变研究所捕捉的关联。

\subsubsection{传统货币变量与行业轮动}
周亮比较广义货币增速、市场利率等指标对货币周期的划分，以考察各周期下的行业配置表现\cite{zhouLiang2018}。游庭玮用申万一级行业面板研究美联储政策利率冲击，提供了按行业特征考察政策敏感度的参照\cite{youTingwei2026}。这一路径关注货币数量和利率等政策变量。本文将报告语言作为单独的特征，在共同样本上比较文本模型与无文本基准，考察文本能否在既有指标之外改善行业排序。

\subsubsection{文本情绪度量与大语言模型}
金融语境会改变通用情绪词的含义，金融语料词典研究已给出直接证据\cite{loughranMcDonald2011}。汪禹、吴彬从中文政策语料扩充并验证倾向词典\cite{wangWu2025}。林建浩、孙乐轩归纳了大语言模型（LLM）在经济金融文本分类、情绪提取和领域词典构建等任务中的应用\cite{linSun2025}。廖海龙将模型生成的情绪因子用于市场择时，余安帮比较了词袋、预训练模型和大语言模型在财务困境预警中的表现\cite{liaoHailong2026,yuAnbang2026}。不同任务的语料与评价标准各异。本文以人工抽检和词典间一致性评估度量，再以同一收益窗口比较指标的预测增量。

\subsubsection{国外央行沟通与资产价格}
Lucca 和 Trebbi 对美联储联邦公开市场委员会声明进行自动化内容分析，考察措辞与利率期限结构的关系\cite{luccaTrebbi2009}。Shapiro 和 Wilson 以美联储内部讨论的情绪衡量决策者对经济形势的判断，并估计其政策目标\cite{shapiroWilson2022}。Ahrens 和 McMahon 从央行演讲提取经济与政策信号\cite{ahrensMcMahon2021}。在直接面向资产定价的研究中，Tetlock 使用财经专栏度量媒体悲观情绪，Ke 等利用财经新闻构造用于预测收益的监督式文本得分\cite{tetlock2007,keKellyXiu2019}。这些研究提供了文本建模与预测评估的参照，各类文本与中国的季度政策报告在受众、信息时效和披露机制上存在区别。

\subsubsection{文献评述与研究定位}
结合项目归档文献与截至 2026 年 11 月 11 日的定向检索，最接近的研究分别覆盖 MPR 与市场整体表现、沟通信号与公司分组反应，以及利率冲击与行业收益。本文聚焦一个不同的证据组合，即使用 MPR 情绪及其与行业特征的交互项，对多个行业的相对收益进行样本外检验。定向检索受数据库和关键词范围约束，本文的贡献以具体研究对象、比较设计和结果证据界定。

\subsection{研究问题与边界}

本文依次检验四个问题。第一，报告情绪能否得到有一致性依据的度量。第二，报告情绪与发布附近的行业异常收益是否关联，以及纳入文本后，跨行业预测表现能否改善。第三，行业融资敏感度、利率敏感度及政策依赖程度能否解释反应差异。第四，LLM 得分能否在词典分数之外提供增量信息。样本覆盖申万一级 31 个行业，实证结论限于行业指数、日与季度收益以及指定模型下的统计关联。报告情绪在同一报告期内对所有行业取相同值，因此行业排序有赖于它与行业特征形成差异化暴露。

\subsection{研究框架与论文结构}

图\ref{fig:framework}展示研究设计。本文先验证数据链路和评分结果，再使用事件研究考察发布附近的反应，随后以横截面、面板和组合检验考察信号在不同目标下的表现。报告期下一季度口径用于呈现原始模型对照。发布后首个完整季度口径用于考察组合持有期表现，并在稳健性分析中按信息可得时点重估预测模型。第\ref{sec:hypotheses}章提出假设，第\ref{sec:data}章说明数据、情绪度量与事件研究，第\ref{sec:models}章报告预测和行业敏感度，第\ref{sec:further}章考察稳健性和误差，第六章总结研究结果与局限。

\paperfigure{fig01_research_framework.png}{研究框架与检验路径（设计阶段）}{fig:framework}{0.99}

\subsection{可能的创新点与贡献}

本文把分析单位从市场整体扩展到申万一级行业指数，检验同一份政策报告是否包含足以区分行业的定价信息。在测量层面，将三套词典、人工抽检及 LLM 得分放在同一语料上比较。在实证层面，联合报告事件窗推断、共同样本的预测比较、组合回测和严格时点复核，明确度量一致性、统计显著性与经济价值之间的区别。实验结果、图表来源和关键处理均有记录，可供后续使用替代词典、文本分段或其他收益窗口时复核。
